\chapter{BMI Calculator}
\index{BMI Calculator}

\section*{Definition and Overview}
A BMI calculator is a tool that works out a person's body mass index (BMI) from their weight and height. BMI is a screening measure used to estimate whether a person is in a healthy weight range, underweight, overweight, or living with obesity. The calculation is simple: weight in kilograms is divided by height in metres squared. BMI calculators are widely available online and in health settings, and they are a useful starting point for discussing weight and health, although they do not measure body fat directly and should be interpreted in the context of the whole person.

\section*{Epidemiology}
BMI is used in health systems around the world, and most adults have had their BMI calculated at some point. Over half the adult population in many developed countries is classified as overweight or obese by BMI, and the proportion has risen steadily over recent decades. BMI is also used in children and adolescents, using age- and sex-specific charts because their bodies are still growing. Although imperfect, BMI remains the most commonly used population measure of weight status because it is cheap, quick, and consistent.

\section*{Aetiology and Pathophysiology}
Understanding how the calculator works helps in interpreting it:

\begin{itemize}
    \item \textbf{The formula:} BMI is weight in kilograms divided by height in metres squared (weight $\div$ height$^2$)
    \item \textbf{Adult categories:} below 18.5 is underweight, 18.5--24.9 is in the healthy range, 25--29.9 is overweight, and 30 and above is obesity, which is itself divided into grades
    \item \textbf{Children and teens:} BMI is compared with age- and sex-specific percentiles rather than fixed adult cut-offs
    \item \textbf{What it reflects:} BMI correlates broadly with body fat but does not measure fat directly, and does not show where fat is carried
    \item \textbf{Ethnic variation:} some population groups, including people of Asian background, are at higher risk of metabolic problems at a lower BMI, so lower cut-offs may be used
\end{itemize}

\section*{Clinical Features}
Using a BMI calculator produces a number and a category, but no physical symptoms. The value of the result lies in what it indicates:

\begin{itemize}
    \item A healthy-range BMI is associated with lower risk of many chronic diseases
    \item A high BMI alerts to increased risk of diabetes, heart disease, stroke, and joint problems
    \item A low BMI may indicate undernutrition, an eating disorder, or another underlying illness
    \item The result should be interpreted alongside waist circumference, muscle mass, and overall health
    \item A single BMI reading is most meaningful when compared over time
\end{itemize}

\section*{Differential Diagnosis}
A BMI result must be interpreted with care:

\begin{itemize}
    \item A muscular person can have a high BMI without excess body fat
    \item Older adults may have a lower BMI but more fat and less muscle than is healthy
    \item Pregnancy, fluid retention, and body frame affect weight
    \item Ethnicity and age change the meaning of a given BMI
    \item BMI alone cannot distinguish who is metabolically healthy or unhealthy
    \item Children should be assessed on growth charts rather than adult cut-offs
\end{itemize}

\section*{When to Seek Medical Care}
Consider seeing a doctor if:

\begin{itemize}
    \item Your BMI is below 18.5 or above 30, or is in the overweight range with other risk factors
    \item You are losing or gaining weight without trying
    \item You have symptoms such as breathlessness, snoring, tiredness, or joint pain
    \item You have diabetes, high blood pressure, or heart disease, or a strong family history
    \item You are concerned about your eating, your body image, or your weight and want support
\end{itemize}

\textbf{Contact your local emergency services for:}
\begin{itemize}
    \item Chest pain, difficulty breathing, or sudden weakness or numbness
    \item Any sudden, severe symptom that worries you
\end{itemize}

\section*{Diagnosis and Investigations}
A single BMI result leads on to a fuller assessment when needed:

\begin{itemize}
    \item \textbf{Correct measurement:} weight on a calibrated scale and height measured accurately
    \item \textbf{Waist circumference:} adds information about abdominal fat, which BMI does not capture
    \item \textbf{Repeat measurements:} tracking BMI over time is more informative than a single reading
    \item \textbf{Blood pressure:} measured as part of a cardiovascular risk check
    \item \textbf{Blood tests:} blood sugar, cholesterol, and liver function may be checked if weight is a concern
    \item \textbf{Clinical assessment:} the doctor considers the whole person, not just the number
\end{itemize}

\section*{Management}
A BMI calculator is a screening tool, not a treatment. Acting on the result involves:

\begin{itemize}
    \item \textbf{Setting goals:} a realistic target, such as modest weight loss of 5--10\% of body weight
    \item \textbf{Diet:} balanced eating, smaller portions, and fewer sugary drinks and highly processed foods
    \item \textbf{Activity:} building up regular physical activity, combining aerobic and strengthening exercise
    \item \textbf{Support:} GP advice, dietitian input, and structured weight-management programs
    \item \textbf{Medical options:} weight-loss medicines or bariatric surgery for people with more severe obesity, under medical supervision
    \item \textbf{For low BMI:} investigation of the cause and nutritional support to gain weight healthily
\end{itemize}

\section*{Complications}
Problems arise not from the calculator itself but from the weight extremes it identifies:

\begin{itemize}
    \item A high BMI increases the risk of diabetes, heart disease, stroke, some cancers, and joint problems
    \item A low BMI increases the risk of infection, osteoporosis, fatigue, and impaired healing
    \item Weight cycling, or repeatedly losing and regaining weight, can be discouraging and hard on health
    \item Misinterpreting BMI can cause anxiety, unhealthy dieting, or dismissal of risk in someone with a normal BMI
    \item Focus on BMI alone may overlook the value of fitness, diet, and mental well-being
\end{itemize}

\section*{Prognosis and Outcomes}
When used correctly, BMI is a helpful guide to weight-related risk and is most useful over time as part of overall health monitoring. People who act on a high BMI with gradual, sustainable changes often see improvements in energy, blood pressure, blood sugar, and quality of life. Likewise, addressing a low BMI with appropriate support restores strength and well-being. The number itself does not determine health -- the habits and follow-up that come after it do.

\section*{Prevention and Self-Care}
\begin{itemize}
    \item Use a BMI calculator correctly, with accurate height and weight measurements
    \item Interpret the result together with waist circumference and your overall health picture
    \item Focus on healthy habits -- balanced meals, regular activity, and good sleep -- rather than on a number alone
    \item Track your BMI occasionally, not obsessively, and look for trends over months
    \item Seek professional advice before starting restrictive diets, especially if your BMI is low
    \item Talk to a doctor about any concerns rather than relying on the calculator as a verdict
\end{itemize}

\section*{Sources and Further Reading}
The following sources provide reliable, up-to-date information on BMI and its interpretation:

\begin{itemize}
    \item World Health Organization. \textit{Body mass index and obesity fact sheets.} https://www.who.int/
    \item NHS. \textit{NHS BMI calculator and healthy weight guidance.} https://www.nhs.uk/conditions/
    \item Centers for Disease Control and Prevention. \textit{About adult BMI and its use.} https://www.cdc.gov/
    \item Mayo Clinic. \textit{Information on BMI and weight management.} https://www.mayoclinic.org/
    \item MedlinePlus. \textit{Body mass index patient information.} https://medlineplus.gov/
    \item NICE. \textit{Obesity identification and management guidelines.} https://www.nice.org.uk/
\end{itemize}
